% Supplementary Material for PAPER5_MONOLITHIC_PIEZO_DRAFT.tex
%   S.1  cubic consistency condition and the OC=SC theorem, in full
%   S.2  elastic-limit gate against the closed-ring solver
%   S.3  finite-element cross-check, provenance and discriminator
%   S.4  e31 inverse: gates, sensitivity, reflection twin
%   S.5  overtones and thickness sweep, full numeric tables
%   S.6  mixed-edge (C-F/F-C) bit-identical reduction proof
%   S.7  force-driven sensing admittance: 12x12 construction and gates
%   S.8  cluster job provenance
\documentclass[11pt]{article}
\usepackage[T1]{fontenc}
\usepackage{lmodern}
\usepackage[margin=1in,letterpaper]{geometry}
\usepackage{amsmath,amssymb,booktabs,array,xcolor}
\usepackage{textcomp}
\usepackage[hidelinks]{hyperref}
\usepackage{graphicx}
\graphicspath{{paper_figures/}{./}{../paper_figures/}}
\usepackage{caption}
\captionsetup{font=small,labelfont=bf,skip=4pt}
\usepackage{placeins}
\usepackage{seqsplit}
\renewcommand{\topfraction}{0.9}
\renewcommand{\bottomfraction}{0.85}
\renewcommand{\textfraction}{0.07}
\renewcommand{\floatpagefraction}{0.85}
\setcounter{topnumber}{4}
\setcounter{bottomnumber}{4}
\setcounter{totalnumber}{8}
\emergencystretch=2.5em
\raggedbottom

\renewcommand{\thesection}{S.\arabic{section}}
\renewcommand{\thesubsection}{S.\arabic{section}.\arabic{subsection}}
\renewcommand{\thetable}{S.\arabic{table}}
\renewcommand{\thefigure}{S.\arabic{figure}}
\renewcommand{\theequation}{S.\arabic{equation}}

\newcommand{\eebar}{\bar{e}_{31}}
\newcommand{\Xibar}{\bar{\Xi}_{33}}
\newcommand{\cbar}{\bar{c}_{11}^{E}}
\newcommand{\fn}[1]{\texttt{\seqsplit{#1}}}

\title{\bfseries Supplementary Material:\\
Exact Radial-Determinant Flexural Frequencies of a Monolithic
Thickness-Poled Piezoelectric Ring: An Open-Circuit-Equals-Short-Circuit Theorem and a Segmented-Electrode Sensing Admittance}
\author{Garrek McCune$^{a,*}$, Daniel Stutts$^{a}$\\
\small\itshape $^a$Mechanical and Aerospace Engineering,\\
Missouri University of Science and Technology, Rolla, MO 65409, USA\\
\small\upshape $^*$Corresponding author. E-mail: ghmkfh@umsystem.edu}
\date{DRAFT --- \today}

\begin{document}
\maketitle
\noindent Main text contains no job numbers; they are listed in
\S\ref{sec:jobs}.

\section{Cubic consistency condition and the open-circuit-equals-short-circuit theorem}
\label{sec:s1}

\subsection{Short-circuit potential and the projected electrical equation}
With both faces grounded, $\phi(\pm H)=0$, write
$\phi=\bar\phi(r,\theta,t)f(z)$ with $f(\pm H)=0$, so
$E_z=-\bar\phi f'$ and $E_\parallel=-f\nabla\bar\phi$. With
$\varepsilon_{rr}+\varepsilon_{\theta\theta}=-z\nabla^2w$, the electric
enthalpy $\mathcal H=\tfrac12\varepsilon\!\cdot\!\bar c\varepsilon
-\eebar E_z(\varepsilon_{rr}+\varepsilon_{\theta\theta})
-\tfrac12\Xibar E_z^2-\tfrac12\Xi_{11}\lvert E_\parallel\rvert^2$
integrates through the thickness to
\begin{equation}
\mathcal H_2=U_{\mathrm{el}}(w)-\eebar I_1\,\bar\phi\,\nabla^2w
-\tfrac12\Xibar I_2\,\bar\phi^2-\tfrac12\Xi_{11}I_3\lvert\nabla\bar\phi\rvert^2,
\end{equation}
with $I_1=\int zf'\,\mathrm dz$, $I_2=\int f'^2\,\mathrm dz$,
$I_3=\int f^2\,\mathrm dz$ over $\lvert z\rvert\le H$. Stationarity in
$w$ and in $\bar\phi$ (Hamilton's principle, the same weight for both
fields) gives the flexural equation (A) below with
$K_{\mathrm{pref}}=-\eebar I_1$, the electrical equation
$\Xi_{11}I_3\Delta\bar\phi-\Xibar I_2\bar\phi+K_{\mathrm{pref}}\Delta
w=0$, and the natural rim condition $\partial\bar\phi/\partial r=0$. In the
long-wave limit $\bar\phi\to(K_{\mathrm{pref}}/\Xibar I_2)\Delta w$ and
$D_{\mathrm{SC}}-d=(\eebar I_1)^2/(\Xibar I_2)$. By Cauchy--Schwarz,
$I_1^2\le I_2\int z^2\,\mathrm dz=\tfrac23H^3I_2$, with equality only for
$f'\propto z$: the quadratic $f=1-z^2/H^2$ is the unique one-term profile
whose long-wave increment is the exact thin-plate value
$\tfrac23H^3\eebar^{\,2}/\Xibar$ (from $\partial_zD_z=0$ with
$\phi(\pm H)=0$). Its integrals are $I_1=-4H/3$, $I_2=8/(3H)$,
$I_3=16H/15$, so
\begin{equation}
M_{rr}^{\mathrm{piezo}}=-\eebar\!\int_{-H}^{H}\!z\,E_z\,\mathrm dz
=-\frac{4H}{3}\,\eebar\,\bar\phi,\qquad
K_{\mathrm{pref}}=\tfrac43H\eebar.
\end{equation}
For comparison, the sinusoidal profile $f=\cos(\pi z/2H)$ gives
$K_{\mathrm{pref}}=(4H/\pi)\eebar$ and a long-wave increment $96/\pi^4$ of
exact under this projection; closing its electrical equation instead by
the thickness-integrated Gauss law (Duan, Quek and Wang's Eq.~(12)
\cite{duan2005}) replaces $\Xibar I_2\bar\phi$ by the face jump of
$D_z$ and gives an increment $12/\pi^2=1.216$ times exact. That closure
was used in an earlier version of this model; its free--free split
($+2.347\%$ against $+1.939\%$) is the $83\%$ finite-element recovery
discussed in \S\ref{sec:s3}. The full moment and shear are
\begin{equation}
M_{rr}=-\Bigl[d\,\partial_r^2 w+(d-2A_1)\bigl(r^{-1}\partial_r w
+r^{-2}\partial_\theta^2 w\bigr)+K_{\mathrm{pref}}\bar\phi\Bigr],\qquad
Q_r=-\Bigl[d\,\partial_r\nabla^2 w+K_{\mathrm{pref}}\partial_r\bar\phi\Bigr],
\end{equation}
with twisting moment $M_{r\theta}=-2A_1\,\partial_r\bigl(r^{-1}\partial_\theta w\bigr)$,
which carries no piezoelectric term. The free-edge natural conditions are
$M_{rr}=0$ and the Kirchhoff effective shear
$Q_r^{\mathrm{eff}}=Q_r+r^{-1}\partial_\theta M_{r\theta}=0$. On a Helmholtz branch
$\Delta Z=\lambda Z$ at angular order $n$ this is
$Q_r^{\mathrm{eff}}[cZ]=-\bigl[KZ'-2A_1n^2(Z'/r^2-Z/r^3)\bigr]$, with
$K=d\lambda+K_{\mathrm{pref}}\chi$. It reduces to $Q_r$ at $n=0$.
The flexural motion equation is
\begin{equation}
d\,\Delta\Delta\bar w+K_{\mathrm{pref}}\,\Delta\bar\phi
-A_2\omega^2\bar w=0. \tag{A}
\end{equation}
The projected electrical equation, divided by $\Xibar I_2$, is
\begin{equation}
a\,\Delta\bar\phi-\bar\phi+c\,\Delta\bar w=0,\qquad
a=\frac{2H^2\Xi_{11}}{5\Xibar},\quad c=\frac{H^2\eebar}{2\Xibar}. \tag{B}
\end{equation}
Positing $\bar\phi=\chi\bar w$ with $\chi$ independent of $r$, (B) forces
$\Delta\bar w=\lambda_B\bar w$ with $\lambda_B=\chi/(a\chi+c)$, and (A)
forces $d\lambda^2+K_{\mathrm{pref}}\chi\lambda-A_2\omega^2=0$.
Eliminating $\lambda$ between the two gives a cubic in $\chi$:
\begin{align}
a_3\chi^3+a_2\chi^2+a_1\chi+a_0&=0,\notag\\
a_3&=K_{\mathrm{pref}}\,a,\quad
a_2=d+K_{\mathrm{pref}}\,c-A_2\omega^2a^2,\notag\\
a_1&=-2A_2\omega^2ac,\quad
a_0=-A_2\omega^2c^2. \label{eq:cubic}
\end{align}
Three roots $\chi_i(\omega)$ give three Helmholtz branches
$\Delta\bar w=\lambda_i\bar w$, two Bessel functions per branch: six
constants, matching the sixth-order coupled system. $\chi=0$ is a root
only at $\eebar=0$ ($a_0\propto\eebar^2\omega^2$), which is why the
purely elastic ring is assembled as a separate $4\times4$
(\S\ref{sec:s2}) rather than by zeroing $\eebar$ inside this cubic.
These coefficients are \emph{not} obtained from the layered plate's own
cubic \cite{mccune2026piezo} by the substitution $h_1\to2H$: the two
potential shapes have opposite parity about their respective plate
midplanes (this ansatz is even in $z$ about the monolithic plate's own
midplane; the layered plate's per-skin interior profile is neither
even nor odd about the sandwich midplane), so the two cubics are genuinely different
operators sharing only their general architecture.

The insulated-rim electrical condition, $\partial\bar\phi/\partial r=0$
at $r_i$ and $r_o$, together with two mechanical rows per edge, fills
the $6\times6$. Free--free uses $M_{rr}=Q_r^{\mathrm{eff}}=0$ at both edges;
clamped--clamped uses $w=w'=0$; clamped--free and free--clamped use one
of each pairing at the two edges (\S\ref{sec:s6}).

\subsection{Open-circuit and the coincidence theorem}
With the bus ansatz $\phi=Vz/(2H)+\bar\phi(1-z^2/H^2)$, the linear
term's field $E_z^{\mathrm{lin}}=-V/(2H)$ is spatially constant --
even in $z$ -- so
$\int_{-H}^{H}z\,\sigma^{\mathrm{lin}}\,\mathrm dz=0$ identically: the
linear piece never enters $M_{rr}$, for any $V$, $H$, or material
constant. Because the bus is a single equipotential surface with no
$r,\theta$ dependence, $E_\parallel^{\mathrm{lin}}=0$, and the enthalpy
cross term $-\Xibar E_z^{\mathrm{lin}}E_z^{f}$ integrates with weight
$\int_{-H}^{H}f'\,\mathrm dz=f(H)-f(-H)=0$; so the linear piece drops from
both (A) and (B): the open-circuit eigenvalue problem for
$(w,\bar\phi)$ is exactly the short-circuit pair, and $V$ is fixed only
afterward, by the scalar charge constraint $Q=0$. The bus charge is the
variable conjugate to $V$ in the same enthalpy,
\begin{equation}
Q=-\frac{\partial}{\partial V}\int_A\mathcal H_2\,\mathrm dA
=-C_0V,\qquad C_0=\frac{\Xibar A}{2H}, \label{eq:Qfull}
\end{equation}
with no motional term: the strain coupling of the linear field carries
the weight $\int_{-H}^{H}z\,\mathrm dz=0$, and $\bar\phi$ does not enter
by the cross-term argument above. Hence $V=0$ on an open-circuit bus for
every flexural mode, and open-circuit coincides with short-circuit.
(Integrating Gauss's law pointwise to the face instead gives a charge
with motional terms in $\int_A\nabla^2w$ and $\int_A\bar\phi$; for the
full face these cancel, the full-face $Q\equiv0$ result of \S\ref{sec:s7},
so the two routes agree on the full electrode.) This is the theorem stated
in the main text: because $V$ never enters $M_{rr}$ -- not merely on a
special mechanical edge, as on the layered plate, but identically --
open-circuit flexure coincides with short-circuit flexure for every
mechanical boundary condition, not only clamped--clamped. On the
layered plate the coincidence is a corollary of $S=0$ under
clamped--clamped conditions specifically \cite{mccune2026piezo}; here it
holds even when $S\ne0$, because the moment channel for $V$ does not
exist at all on this stacking.

\section{Elastic-limit gate}
\label{sec:s2}
Setting $\eebar=0$, $\bar\phi=0$ in (A) recovers the isotropic Kirchhoff
ring equation with $D=d=\tfrac23\cbar H^3$,
$\nu=\nu_p=\bar c_{12}^E/\cbar$, plate thickness $h=2H$. Matching the
standard isotropic bending stiffness $D=Eh^3/[12(1-\nu^2)]$ requires
$E=\cbar(1-\nu_p^2)$. At the material point of the main text this gives
$\cbar=85.661\,\mathrm{GPa}$, $\bar c_{12}^E=24.661\,\mathrm{GPa}$,
$\nu_p=0.28789$, $E=78.561\,\mathrm{GPa}$. The free--free $n=0$
fundamental of the new mixed-edge $4\times4$ at these constants agrees
with an independently implemented closed-annulus Kirchhoff solver
\cite{mccune2026ring} at $4.77\times10^{-10}$ relative. The coupled
$6\times6$ is not exercised at this gate (\S\ref{sec:s1}: $\chi=0$ is a
root only at $\eebar=0$, so the coupled operator is the wrong tool at
exact zero coupling). This gate passed before any coupled number in the
main text was trusted, and uses no piezoelectric code path.

\section{Finite-element cross-check}
\label{sec:s3}
Axisymmetric coupled-field finite elements (piezoelectric-plane
elements, KEYOPT(1)=1001, Lanczos eigensolver), millimetre--newton--
second--tonne units with relative permittivity, one PZT-4 rectangle
$\lvert z\rvert\le H=10\,\mathrm{mm}$ (no steel host), cloned from the
mesh, material block, and electrical-boundary conventions already
confirmed for the layered plate's own free--free family
\cite{mccune2026piezo}, with the host volume deleted and the piezo
volume filling the full thickness. Three decks on the same mesh:
coupling-off (\texttt{e0}), both-faces short-circuit (\texttt{scboth},
VOLT$=0$ at $z=\pm H$), and open-circuit (\texttt{oc}, VOLT$=0$ at
neither face, a single coupled VOLT set on both faces together,
unconstrained). Job \textbf{2520476}, \texttt{compute-33-02.mill},
$3/3$ \texttt{RUN COMPLETED}, \textbf{0 errors, 0 warnings} on every
deck, scored from the raw transcript and mode/discriminator files.

Mode~1 is the free--free rigid body ($0\,\mathrm{Hz}$) on every deck.
Mode~2 is the first bending mode, identified by the through-thickness
discriminator at $r=r_o$: transverse displacement same sign at
$z=\pm H$, in-plane displacement opposite sign, ratio $\approx19.7$ on
all three decks.

\begin{table}[ht]
\centering
\caption{F--F $n=0$ finite-element mode~2 (bending), job 2520476.}
\label{tab:fe}
\small
\begin{tabular}{@{}lccc@{}}
\toprule
deck & frequency [Hz] & closed form [Hz] & rel.\ vs closed form \\
\midrule
\texttt{e0} (coupling off) & 73.506288 & 73.593 & $-0.118\%$ \\
\texttt{scboth} (SC) & 74.933783 & 75.020 & $-0.115\%$ \\
\texttt{oc} (OC) & 74.933783 & 75.020 & $-0.115\%$ \\
\bottomrule
\end{tabular}
\end{table}

\texttt{scboth} and \texttt{oc} mode~2 are bit-identical at the printed
eight significant figures, $74.933783\,\mathrm{Hz}$ on both -- the
coincidence theorem confirmed in three-dimensional coupled-field finite
elements, not only in the Kirchhoff $6\times6$ -- and so are the next
several flexural overtones on the same pair of decks. \texttt{scboth}
sits $1.942\%$ above \texttt{e0}, against a closed-form $+1.939\%$
(ratio $1.001$); the finite-element short-circuit frequency sits
$-0.115\%$ from the closed form, the same offset as the coupling-off
deck ($-0.118\%$), so the residual is the elastic Kirchhoff-versus-2D
difference, not the coupling. \textbf{History.} The first version of
this model closed the electrical equation by the thickness-integrated
Gauss law with a sinusoidal profile (\S\ref{sec:s1}) and predicted
$+2.347\%$ here; finite elements recovered $83\%$ of it. The same
closure predicted $+0.961\%$ for a clamped PZT-4 annular sector
($r_i/r_o=0.03/0.07$, $2H=0.8\,\mathrm{mm}$, $2\Theta=\pi$), where
three-dimensional SOLID226 finite elements give $+0.782\%$ (ratio
$0.814$); the consistent projection gives $+0.796\%$ for that sector.
Both ratios are $\pi^2/12=0.822$ to within the elastic offsets, which is
how the projection was identified. Neither comparison was used to set a
coefficient: the quadratic and its projection are fixed by the
derivation of \S\ref{sec:s1}.

A companion extensional mode on the same mesh (mode~6 on \texttt{e0}
and \texttt{scboth}) is $1669.57\,\mathrm{Hz}$ on both, and
$1699.58\,\mathrm{Hz}$ on \texttt{oc} ($+1.80\%$): the linear
(membrane) piezoelectric field does couple to extension, confirming
the negative control of main-text \S2.1 in finite elements as well as
in closed form. This is disclosed as claim~C of this series
\cite{mccune2026piezo} and is not built further here.

\section{$e_{31}$ inverse: gates, sensitivity, reflection twin}
\label{sec:s4}
With $\cbar$ held at its known value, $e_{31}$ is inverted from the
free--free (\texttt{coupled\_bisect}) and clamped--clamped
(\texttt{cc\_coupled\_bisect}) short-circuit fundamentals of
Table~1 of the main text. Job \textbf{2520482},
\texttt{compute-33-02.mill}, 562.0~s, dps$=60$/inner\_iters$=45$/
outer\_iters$=30$, \texttt{.err} empty (0 bytes).

\textbf{G\_worker.} The worker-process scalar reconstruction reproduces
the in-process bisection roots at rel$=0.0$ on both legs, before any
inversion built on the in-process class is trusted.

\textbf{G0 (round trip).} Bracket $[1.0,8.0]\,\mathrm{C\,m^{-2}}$ on the
physical branch (see the reflection-twin derivation below for why this
range, not $[1,10]$). Recovered $\hat e_{31}=4.100000001$ against true
$4.1$, rel\_err $=2.499\times10^{-10}$ (bar $10^{-6}$), identical on
both legs.

\textbf{G1 (negative controls).} A wrong-sign bracket
$[-8.0,-1.0]$ shows no sign flip; $e_{31}=-4.1$ gives a frequency
$>5\%$ from target on both legs; $e_{31}=0$ raises the singular-cubic
exception of \S\ref{sec:s1} when routed through a real determinant call
(not at construction alone).

\textbf{G2 (sensitivity).} F--F:
$\mathrm d\omega/\mathrm de_{31}=-3.617563\,\mathrm{rad\,s^{-1}}$ per
$\mathrm{C\,m^{-2}}$; propagated uncertainty $0.3178\%$ of $e_{31}$ at
$0.01\%$ frequency precision, $3.1780\%$ at $0.1\%$. C--C:
$-12.59036\,\mathrm{rad\,s^{-1}}$ per $\mathrm{C\,m^{-2}}$; $0.3405\%$
and $3.4050\%$. Both legs agree closely with the independent long-wave
closed-form estimate of main-text \S2.4 ($\sim0.3\%$ of $e_{31}$ at
$0.01\%$ frequency if stiffness is known). (The Gauss-law closure of
\S\ref{sec:s1} gave $0.2649\%$ and $0.2822\%$; its larger split made
$e_{31}$ look better conditioned, and fitted to a measured or
finite-element split it would return $\lvert\eebar\rvert$ low by
$(\pi^2/12)^{1/2}=0.907$.)

\textbf{Reflection twin, derived.} Let
$\eebar=e_{31}-(c_{13}^E/c_{33}^E)e_{33}$. The cubic coefficients of
\eqref{eq:cubic} are odd in $\eebar$ for $a_3,a_1$ and even for
$a_2,a_0$ ($K_{\mathrm{pref}}$ and $c$ are odd, $a$ is even). Flipping $\eebar\to-\eebar$ therefore maps the root set
$\{\chi_i\}\to\{-\chi_i\}$ while each $\lambda_i$ is invariant (numerator
and denominator of $\lambda_B(\chi)$ flip together) and
$K_{\mathrm{pref}}\chi$ is invariant termwise ($K_{\mathrm{pref}}$ and
$\chi$ each flip, the product does not): the determinant's $M$-row and
$Q$-row blocks are unchanged, and only the $\bar\phi'$-row block flips
sign as a whole -- an overall row rescaling that does not move where the
determinant vanishes. Hence
$e_{31}^{\mathrm{twin}}=2(c_{13}^E/c_{33}^E)e_{33}-e_{31}$ gives a
bit-identical root frequency to $e_{31}$ itself. At $e_{31}=4.1$,
$e_{31}^{\mathrm{twin}}=13.800869565217392$, confirmed to rel$=0$ on
both legs. This is why the inversion bracket is $[1.0,8.0]$ and not a
wider $[1,10]$: a wider bracket would cross $\eebar=0$
($e_{31}\approx8.9504\,\mathrm{C\,m^{-2}}$, where the cubic degenerates
in degree) and approach the twin itself. Do not chase the twin as a
second physical solution; it is the same root, not a different one.

\section{Overtones and thickness sweep}
\label{sec:s5}
\begin{sloppypar}
A sign-flip scan (dps$=60$/inner\_iters$=45$, production settings;
job \textbf{2521175}, \texttt{compute-34-07.mill}, 421.8~s,
\texttt{.err} empty) found three F--F pairs in
$[455,6000]\,\mathrm{rad\,s^{-1}}$ and two C--C pairs in
$[1700,6000]$; every pair reproduces the reflection twin of
\S\ref{sec:s4} at $\sim5$--$6\times10^{-16}$ relative, not only at the
fundamental.
\end{sloppypar}

\begin{table}[ht]
\centering
\caption{F--F and C--C $n=0$ overtones. Elastic and short-circuit
($=$ open-circuit) roots in $\mathrm{rad\,s^{-1}}$.}
\label{tab:overtones}
\small
\begin{tabular}{@{}llccc@{}}
\toprule
edge & mode & elastic & SC $=$ OC & split \\
\midrule
F--F & fundamental & 462.398 & 471.365 & $+1.939\%$ \\
F--F & OT1 & 2169.893 & 2203.925 & $+1.568\%$ \\
F--F & OT2 & 5270.649 & 5355.495 & $+1.610\%$ \\
C--C & fundamental & 1726.890 & 1757.694 & $+1.784\%$ \\
C--C & OT1 & 4776.965 & 4861.470 & $+1.769\%$ \\
\bottomrule
\end{tabular}
\end{table}

Short-circuit sits above elastic at every pair, both legs -- the
main-text stiffening-sign theorem generalizes beyond the fundamental.
The split trend is not monotonic: F--F dips from $1.94\%$ to $1.57\%$
at the first overtone, then ticks back up to $1.61\%$ at the second;
C--C stays flatter, $1.78\%\to1.77\%$.

A nine-point thickness sweep (F--F $n=0$ only, fixed $r_i,r_o$, job
\textbf{2521176}, same host, concurrent, 805.8~s, \texttt{.err}
empty) tests the long-wave prediction that the split is independent of
$H$ at leading order.

\begin{table}[ht]
\centering
\caption{F--F $n=0$ fundamental versus ceramic thickness, fixed radii.
Long-wave closed-form estimate of the thin-limit split: $1.814\%$.}
\label{tab:thickness}
\small
\begin{tabular}{@{}lccc@{}}
\toprule
$H/r_o$ & elastic [rad/s] & SC [rad/s] & split \\
\midrule
0.005 & 138.719 & 141.410 & $1.9399\%$ \\
0.010 & 277.439 & 282.820 & $1.9397\%$ \\
0.01667 & 462.398 & 471.365 & $1.9392\%$ \\
0.025 & 693.596 & 707.040 & $1.9383\%$ \\
0.040 & 1109.754 & 1131.234 & $1.9355\%$ \\
0.060 & 1664.632 & 1696.760 & $1.9300\%$ \\
0.090 & 2496.947 & 2544.843 & $1.9182\%$ \\
0.130 & 3606.702 & 3675.111 & $1.8967\%$ \\
0.180 & 4993.895 & 5086.967 & $1.8637\%$ \\
\bottomrule
\end{tabular}
\end{table}

Every point returns exactly one elastic and one coupled root in its
search window; short-circuit exceeds elastic at every point; the
thinnest point's split ($1.9399\%$) is within a factor of $1.069$ of the
raw long-wave estimate $\tfrac12\eebar^{\,2}/(\Xibar\cbar)=1.8141\%$
(unrounded constants; the ratio exceeds one because the same
$\eebar^{\,2}/\Xibar$ also stiffens $\bar c_{12}^E$, which the
bending-only estimate omits). The
$H$-independence prediction holds well below $H/r_o\approx0.025$
(split flat near $1.94\%$) and then drifts down to $1.8637\%$ by
$H/r_o=0.18$, a $-3.93\%$ relative drift from the thinnest point -- a
real, small finite-thickness correction beyond the long-wave leading
order, consistent with Kirchhoff kinematics being best justified in
exactly the thin regime the main-text geometry ($H/r_o\approx0.0167$)
already sits in.

\section{Mixed-edge (C--F/F--C) bit-identical reduction proof}
\label{sec:s6}
Clamped--free and free--clamped determinants are built from a single
mixed-edge row-selection routine that picks, at each edge
independently, either the $(w,w')$ pair (clamped) or the $(M_{rr},Q_r^{\mathrm{eff}})$
pair (free), following the pattern already used for the elastic ring's
own mixed-edge family \cite{mccune2026ring}, adapted to the three-branch
cubic of \S\ref{sec:s1}. Because free--free and clamped--clamped are
literal special cases of this routine (both edges set to the same
label), the reduction was checked bit-for-bit (\texttt{mpmath} value
equality, not a tolerance), not only re-derived: at $n=0,\ \omega=500$
and $n=2,\ \omega=1500$, the mixed-edge routine with both edges 'F'
(respectively 'C') reproduces the original free--free (respectively
clamped--clamped) elastic and coupled determinants exactly, eight of
eight comparisons. This is a stronger correctness guarantee than an
independent derivation would give: a defect in the row-selection logic
would almost certainly have broken the already-validated free--free and
clamped--clamped paths too, not only left a silent error in the new
clamped--free/free--clamped paths.

Physical anchors (job \textbf{2521220}, \texttt{compute-34-04.mill},
369.0~s, \texttt{.err} empty): clamped--free
elastic$=1605.4427637454191$, SC$=1637.9221255770667$
($+2.023\%$); free--clamped elastic$=556.8844158380273$,
SC$=564.9003213461377$ ($+1.439\%$). Both stiffen, consistent with
the main text's stiffening-sign estimate (Eq.~7). The negative control ($e_{31}=0$
refused) still fires through the new coupled methods. A worker-path
regression -- new scalar-reconstruction functions for each of the four
new bisection methods -- was confirmed bit-identical to in-process, both
in a standalone check and in the local test suite, for all four new
methods.

The genuine physical finding of main-text \S4.2 -- the inner edge
dominating bending stiffness far more than the outer edge at this
$r_i/r_o=1/6$ -- is reported as an empirical fact of this one geometry,
with a plausible but unproven mechanism: at $n=0$ the free-edge moment
condition's $r$-dependent term is $-(2A_1/r)\,Z'$, which is
$r_o/r_i=6\times$ more sensitive to a clamp override at $r_i$ than at
$r_o$, consistent with, but not a rigorous proof of, the observed
asymmetry (clamped--free at $90.4\%$ of the way from free--free to
clamped--clamped; free--clamped at $7.5\%$).

\section{Force-driven sensing admittance: construction and gates}
\label{sec:s7}
The two-region twelve-unknown system of the main text's numerical-method
section is
built from the same three $\chi$-cubic branches as \S\ref{sec:s1}, split
at the load radius $r_F$, with column order (region~I then region~II,
three branches, two Bessel functions each) matching the free-vibration
determinant's own inner loop. On a Helmholtz function $\Delta Z=\lambda
Z$, the physical $n=0$ resultants collapse to
$M_{rr}[cZ]=-(KZ-2A_1Z'/r)$, $Q_r[cZ]=-KZ'$,
$K=d\lambda+K_{\mathrm{pref}}\chi$. Six rim rows (three per edge, F--F
SC) plus six interface rows -- $[w]=[w']=[M_{rr}]=[\bar\phi]
=[\bar\phi']=0$ and the one inhomogeneous shear jump
$[Q_r]=-F/(2\pi r_F)$ -- close the twelve-unknown system, solved by an
equilibrated linear solve (row and column scaling, tracked and
un-scaled after solving) rather than a determinant search, because the
third (evanescent) $\chi$-branch produces radial arguments large enough
that raw entries span more than fifty orders of magnitude at this
geometry.

\begin{sloppypar}
Five gates, all in-sandbox and independently cluster-confirmed
(job \textbf{2521701}, \texttt{compute-34-04.mill}, 247.5~s,
\texttt{.err} empty, every printed and logged number bit-identical
between sandbox and cluster):
\end{sloppypar}

\begin{itemize}
\item \textbf{G1, force-balance identity} (the load-bearing sign gate on
the shear-jump right-hand side):
$2\pi\int_{r_i}^{r_o}r\,w(r)\,\mathrm dr=-F/(A_2\omega^2)$ at
rel\_err $=1.112761845\ldots\times10^{-61}$ -- essentially exact,
confirming the jump's sign without assuming it.
\item \textbf{G2, pole location} (unique-continuation gate): the driven
system's own homogeneous determinant and the free-vibration
short-circuit bisection agree at the F--F fundamental,
$471.3646021920983\,\mathrm{rad\,s^{-1}}$, rel\_err $=0.0$ on every
computed digit.
\item \textbf{G3, $\eebar\to0$ limit}, scaled correctly (raw $e_{31}$
and $e_{33}$ together, so $\eebar$ itself scales to zero -- see the note
below): $Q_{\mathrm{in}}$ at scale $1,10^{-1},\ldots,10^{-4}$ is
$5.004\times10^{-12}$, $6.627\times10^{-13}$, $6.648\times10^{-14}$,
$6.648\times10^{-15}$, $6.648\times10^{-16}$; decade ratios $7.551$,
$9.968$, $10.000$, $10.000$; log--log slope over the three deep decades
is $0.999993$, clean linear vanishing.
\item \textbf{G\_mu} (diagnostic): the two mechanical Helmholtz branches
match the purely elastic wavenumber $\pm\mu=\pm10.250145149520318$ at
rel\_err $=0.0$ at scale $10^{-8}$, and the third branch converges to
the pure-electrical Helmholtz eigenvalue $1/a=5\Xi_{33}/(2H^2\Xi_{11})
=20497.61$ of (B) (with $\Xibar\to\Xi_{33}$, because $e_{33}$ is scaled
to zero together with $e_{31}$).
\item \textbf{G\_worker}: the worker-process point evaluation reproduces
$Y_{\mathrm{sense}}$ bit-identically (real and imaginary parts) both at
$\omega=200$ ($Y_{\mathrm{sense}}=5.00405914569838\times10^{-12}$) and
at $\omega=473$, just above the pole
($Y_{\mathrm{sense}}=6.062573639651388\times10^{-10}$).
\end{itemize}

\textbf{Methodological note.} An earlier pass at G3, scaling only the
raw constructor keyword \texttt{e31} toward zero while leaving
\texttt{e33} fixed, appeared to show $Q_{\mathrm{in}}$ plateauing near
$1.04\times10^{-11}$ instead of vanishing, and the two mechanical
branches plateauing near $\lambda\approx\pm9.557$ instead of the
correct $\pm10.250145\ldots$. Root cause: $\eebar=e_{31}
-(c_{13}^E/c_{33}^E)e_{33}$ is dominated by the unscaled \texttt{e33}
cross term at this material point, so with \texttt{e33} held fixed,
$\eebar$ plateaus near $-8.951$ (its value at raw $e_{31}=0$) and never
actually approaches zero -- the singular limit being probed was never
reached. This was root-caused by direct numerical probe, consistent with
this series' own methodology: precision was ruled out first (\texttt{dps}
$50/100/150$ gave bit-identical wrong answers under the buggy scaling,
a real mathematical limit of the wrong quantity, not a conditioning
artifact), which is what forced tracing the discrepancy to the test
harness rather than the solver. A permanent regression test pins the
correct behavior down going forward.

Reported, not gated: sweeping $Y_{\mathrm{sense}}(\omega)$ across the
F--F short-circuit fundamental shows the resonance of main-text
Figure~3 -- magnitude climbing from $\sim10^{-10}$ well off-resonance to
$\pm9.76\times10^{-8}$ within $0.01\,\mathrm{rad\,s^{-1}}$ of
$471.3646$, with a clean sign flip through the pole, textbook
driven-response behavior for a transfer function with a simple pole at
the mechanical resonance and no damping in the operator.

\section{Cluster job provenance}
\label{sec:jobs}
All jobs on The Mill \cite{mill2024}, solver snapshot
\texttt{2026-07-10.s10}. The monolithic piezoelectric path is a new
module, unused by the elastic or layered-piezoelectric defaults.

\begin{itemize}
\item 2520476 --- F--F finite-element queue (\texttt{e0}/\texttt{scboth}/
      \texttt{oc}), $3/3$ \texttt{RUN COMPLETED}, 0 errors, 0 warnings.
      Mode~2 bit-identical OC/SC at $74.933783\,\mathrm{Hz}$.
\item 2520482 --- F--F/C--C $e_{31}$ inverse, \texttt{SENTINEL
      PASS\_ALL}, both legs, $\hat e_{31}$ rel\_err $2.499\times10^{-10}$.
\item 2521175 --- $n=0$ overtones (F--F/C--C), \texttt{SENTINEL
      PASS\_ALL}, bit-identical to the in-sandbox smoke-test digits.
\item 2521176 --- thickness sweep (F--F, 9 points), \texttt{SENTINEL
      PASS\_ALL}, \textbf{G\_found}, \textbf{G\_theorem}, and
      \textbf{G\_longwave} all pass.
\item 2521220 --- C--F/F--C mixed-edge, \texttt{SENTINEL PASS\_ALL},
      \textbf{G\_reduce} bit-identical F--F/C--C reduction, worker-path green.
\item 2521701 --- force-driven sensing admittance, \texttt{SENTINEL
      PASS\_ALL}, all 5 gates, bit-identical to the in-sandbox run.
\end{itemize}
\noindent\textbf{Consistent-projection rerun (2026-09-23).} Jobs
2520482, 2521175, 2521176, 2521220 and 2521701 used the Gauss-law
closure of \S\ref{sec:s1}. All five probes were re-run unchanged
against the corrected solver (default \texttt{projection='consistent'};
the old closure is kept as \texttt{projection='duan'} and reproduces the
old numbers bit-for-bit). Every gate passes (\texttt{SENTINEL PASS\_ALL}
on all five), and every number in this document and the main text is
from that rerun. The rerun was done in-sandbox at the probes' production precision; it was not repeated on the cluster, because every earlier cluster archival run of these same probes reproduced its in-sandbox run bit-for-bit. The finite-element job 2520476 is unchanged: it does not
depend on the analytic model.

\clearpage
\begin{thebibliography}{99}\small
\setlength{\itemsep}{0pt}
\setlength{\parsep}{0pt}
\bibitem{duan2005} W.H. Duan, S.T. Quek, Q. Wang, ``Free vibration analysis
of piezoelectric coupled thin and thick annular plate,'' J. Sound Vib. 281
(2005) 119--139.
\bibitem{mccune2026ring} G. McCune, D. Stutts, ``Exact radial-determinant
frequencies of annular and circular plates: the closed-ring and solid-disk
limits of the sector wave-function method,'' submitted (2026).
\bibitem{mccune2026piezo} G. McCune, D. Stutts, ``Exact radial-determinant
frequencies of a piezoelectric annular plate, and inverse recovery of
$e_{31}$ from open-circuit flexure,'' submitted (2026).
\bibitem{mill2024} Missouri University of Science and Technology, ``The Mill
HPC Cluster,'' \emph{The Mill} (2024), Article 1.
\url{https://scholarsmine.mst.edu/the-mill/1}.
\end{thebibliography}
\end{document}
